\documentclass[11pt]{article}

\usepackage[utf8]{inputenc}
\usepackage[T1]{fontenc}
\usepackage{amsmath, amssymb, amsthm}
\usepackage{graphicx}
\usepackage{booktabs}
\usepackage[round]{natbib}
\usepackage[hidelinks]{hyperref}

\graphicspath{{figures/}}

\theoremstyle{plain}
\newtheorem{theorem}{Theorem}
\newtheorem{lemma}{Lemma}
\theoremstyle{definition}
\newtheorem{definition}{Definition}

\title{A Title for the Paper}
\author{%
  First Author\thanks{Affiliation, \texttt{first@example.com}}
  \and
  Second Author\thanks{Affiliation, \texttt{second@example.com}}
}
\date{\today}

\begin{document}

\maketitle

\begin{abstract}
A concise summary of the problem, the approach, and the main result.
\end{abstract}

\section{Introduction}
\label{sec:intro}

Motivation, contribution, and an outline of the paper. Cite related work
like \citet{example2020} or \citep{example2020}.

\section{Method}
\label{sec:method}

\begin{definition}
State definitions here.
\end{definition}

\begin{theorem}
State the main result here.
\end{theorem}

\begin{proof}
Sketch or full proof.
\end{proof}

\section{Results}
\label{sec:results}

% \begin{figure}[t]
%   \centering
%   \includegraphics[width=0.7\linewidth]{example}
%   \caption{Caption.}
%   \label{fig:example}
% \end{figure}

\section{Conclusion}
\label{sec:conclusion}

Summarize and point to future work.

\bibliographystyle{plainnat}
\bibliography{references}

\end{document}
